\chapter{DL-016 Mode-aware runtime theme adapter}\label{record-181}

\noindent{\footnotesize\color{BookMuted}\path{docs/design/DL-016-mode-aware-runtime-theme-adapter.md}\par}\smallskip

\begin{RecordMeta}
title: DL-016 Mode-aware runtime theme adapter
legend: DL
lane: bad-code
priority: medium
github_issue: 313
keywords:
  - design-language
  - themes
  - modes
  - app-frame
  - dogfood

\end{RecordMeta}

Legend: \href{https://github.com/flyingrobots/bijou/blob/1853ad902f6bf219326409e7f743a5a1d6517dbc/docs/legends/DL-design-language.md}{DL - Design Language}

\section{Decision Summary}\label{record-181-decision-summary}

Bijou should distinguish theme identity from theme mode at the app-frame boundary.

The existing runtime \texttt{T\allowbreak{}h\allowbreak{}e\allowbreak{}m\allowbreak{}e} object is a concrete palette. That stays true for compatibility. The new shape is an adapter around that concrete palette:

% Presentation of original arrow notation. See flow-manifest.json.
\begin{center}
\begin{tikzpicture}[book flow node/.style={draw=BookHeading,line width=.65pt,fill=BookPanel,rounded corners=1.5pt,align=center,inner sep=5pt,minimum height=10mm,font=\fontsize{9}{11}\selectfont,execute at begin node={\hyphenpenalty=10000\exhyphenpenalty=10000}},book flow edge/.style={-{Latex[length=2mm]},draw=BookHeading,line width=.7pt}]
\node[book flow node,text width=58mm] (n0) at (0,0.0) {Frame\allowbreak{}Shell\allowbreak{}Theme\allowbreak{}Family, id: "dogfood"};
\node[book flow node,text width=58mm] (n1) at (4,0.0) {dark: Theme("dogfood-dark")};
\node[book flow node,text width=58mm] (n2) at (4,-1.8) {light: Theme("dogfood-light")};
\draw[book flow edge] (n0) -- (n1);
\draw[book flow edge] (n0) -- (n2);
\end{tikzpicture}
\end{center}


Single-mode shell themes remain valid:

\begin{CodeBlock}
FrameShellTheme
  id: "default"
  theme: Theme("default")
\end{CodeBlock}

The frame resolves both forms into concrete choices before rendering. A choice has a stable id, a family id, an optional mode id, and a resolved concrete theme. This gives settings, inspectors, and DOGFOOD one product identity while keeping the legacy render path stable.

\section{Linked Work}\label{record-181-linked-work}

\begin{itemize}
\tightlist
\item
  User story: \href{https://github.com/flyingrobots/bijou/issues/313}{\#313}
\item
  Token builder foundation: \href{https://github.com/flyingrobots/bijou/issues/308}{\#308}
\item
  Theme Inspector follow-up: \href{https://github.com/flyingrobots/bijou/issues/311}{\#311}
\item
  Theme Lab follow-up: \href{https://github.com/flyingrobots/bijou/issues/315}{\#315}
\end{itemize}

\section{Sponsored Perspectives}\label{record-181-sponsored-perspectives}

\subsection{Human Perspective}\label{record-181-human-perspective}

A designer-engineer wants DOGFOOD to present one DogFood theme with dark and light modes so that switching color scheme does not look like switching product identity.

\subsection{Agent Perspective}\label{record-181-agent-perspective}

An implementation agent wants app-frame theme facts to carry both family and mode identity so tests, inspectors, and future provenance tools can prove which semantic theme and concrete mode produced the current cells.

\section{Hill}\label{record-181-hill}

An app can declare a mode-aware shell theme, switch between its dark and light modes through the existing settings drawer, and observe the selected family and mode without breaking consumers that still pass single concrete themes.

\section{Problem}\label{record-181-problem}

DL-013 made theme authoring mode-aware: one theme owns dark and light token values. The runtime still consumes one concrete \texttt{T\allowbreak{}h\allowbreak{}e\allowbreak{}m\allowbreak{}e} at a time. PR \#310 therefore had to model DOGFOOD as two independent shell themes:

\begin{CodeBlock}
dogfood-dark
dogfood-light
\end{CodeBlock}

That representation works mechanically, but it teaches the wrong model:

\begin{itemize}
\tightlist
\item
  theme identity and mode identity are conflated
\item
  settings rows mix product themes and color modes in one flat list
\item
  Theme Inspector and Theme Lab cannot report a stable family id
\item
  tests assert \texttt{d\allowbreak{}o\allowbreak{}g\allowbreak{}f\allowbreak{}o\allowbreak{}o\allowbreak{}d\allowbreak{}-\allowbreak{}d\allowbreak{}a\allowbreak{}r\allowbreak{}k} and \texttt{d\allowbreak{}o\allowbreak{}g\allowbreak{}f\allowbreak{}o\allowbreak{}o\allowbreak{}d\allowbreak{}-\allowbreak{}l\allowbreak{}i\allowbreak{}g\allowbreak{}h\allowbreak{}t} as peer themes instead of \texttt{d\allowbreak{}o\allowbreak{}g\allowbreak{}f\allowbreak{}o\allowbreak{}o\allowbreak{}d} plus a selected mode
\end{itemize}

\section{Scope}\label{record-181-scope}

This slice adds the first runtime adapter, not a full replacement for \texttt{T\allowbreak{}h\allowbreak{}e\allowbreak{}m\allowbreak{}e}.

\begin{itemize}
\tightlist
\item
  Add a mode-aware \texttt{F\allowbreak{}r\allowbreak{}a\allowbreak{}m\allowbreak{}e\allowbreak{}S\allowbreak{}h\allowbreak{}e\allowbreak{}l\allowbreak{}l\allowbreak{}T\allowbreak{}h\allowbreak{}e\allowbreak{}m\allowbreak{}e\allowbreak{}F\allowbreak{}a\allowbreak{}m\allowbreak{}i\allowbreak{}l\allowbreak{}y} path beside concrete \texttt{F\allowbreak{}r\allowbreak{}a\allowbreak{}m\allowbreak{}e\allowbreak{}S\allowbreak{}h\allowbreak{}e\allowbreak{}l\allowbreak{}l\allowbreak{}T\allowbreak{}h\allowbreak{}e\allowbreak{}m\allowbreak{}e}.
\item
  Resolve shell theme specs into flattened runtime choices.
\item
  Preserve old single-theme specs unchanged.
\item
  Include selected mode facts in \texttt{F\allowbreak{}r\allowbreak{}a\allowbreak{}m\allowbreak{}e\allowbreak{}S\allowbreak{}h\allowbreak{}e\allowbreak{}l\allowbreak{}l\allowbreak{}T\allowbreak{}h\allowbreak{}e\allowbreak{}m\allowbreak{}e\allowbreak{}C\allowbreak{}h\allowbreak{}a\allowbreak{}n\allowbreak{}g\allowbreak{}e}.
\item
  Convert DOGFOOD to one \texttt{d\allowbreak{}o\allowbreak{}g\allowbreak{}f\allowbreak{}o\allowbreak{}o\allowbreak{}d} shell theme with \texttt{d\allowbreak{}a\allowbreak{}r\allowbreak{}k} and \texttt{l\allowbreak{}i\allowbreak{}g\allowbreak{}h\allowbreak{}t} modes.
\item
  Update Theme Inspector and Theme Lab facts to use family and mode labels.
\item
  Add tests proving mode switching, compatibility, and context isolation.
\end{itemize}

\section{Non-Goals}\label{record-181-non-goals}

\begin{itemize}
\tightlist
\item
  Do not replace the concrete \texttt{T\allowbreak{}h\allowbreak{}e\allowbreak{}m\allowbreak{}e} interface.
\item
  Do not require every app to adopt mode-aware themes.
\item
  Do not add terminal background-color detection here.
\item
  Do not add element-level token provenance.
\item
  Do not implement Theme Inspector mode preview/apply UX from \#311 in this slice.
\end{itemize}

\section{Runtime Contract}\label{record-181-runtime-contract}

\texttt{F\allowbreak{}r\allowbreak{}a\allowbreak{}m\allowbreak{}e\allowbreak{}S\allowbreak{}h\allowbreak{}e\allowbreak{}l\allowbreak{}l\allowbreak{}T\allowbreak{}h\allowbreak{}e\allowbreak{}m\allowbreak{}e\allowbreak{}S\allowbreak{}p\allowbreak{}e\allowbreak{}c} accepts either a concrete \texttt{F\allowbreak{}r\allowbreak{}a\allowbreak{}m\allowbreak{}e\allowbreak{}S\allowbreak{}h\allowbreak{}e\allowbreak{}l\allowbreak{}l\allowbreak{}T\allowbreak{}h\allowbreak{}e\allowbreak{}m\allowbreak{}e} with \texttt{t\allowbreak{}h\allowbreak{}e\allowbreak{}m\allowbreak{}e} or a mode-aware \texttt{F\allowbreak{}r\allowbreak{}a\allowbreak{}m\allowbreak{}e\allowbreak{}S\allowbreak{}h\allowbreak{}e\allowbreak{}l\allowbreak{}l\allowbreak{}T\allowbreak{}h\allowbreak{}e\allowbreak{}m\allowbreak{}e\allowbreak{}F\allowbreak{}a\allowbreak{}m\allowbreak{}i\allowbreak{}l\allowbreak{}y} with a non-empty \texttt{m\allowbreak{}o\allowbreak{}d\allowbreak{}e\allowbreak{}s} array.

\begin{CodeBlock}
const theme = {
  id: 'dogfood',
  label: 'DOGFOOD',
  modes: [
    { id: 'dark', label: 'Dark', theme: dogfoodDark },
    { id: 'light', label: 'Light', theme: dogfoodLight },
  ],
} satisfies FrameShellThemeSpec;
\end{CodeBlock}

A mode-aware entry resolves into choices:

\begin{CodeBlock}
choice id: dogfood:dark
family id: dogfood
mode id: dark
label: DOGFOOD / Dark
theme: dogfood-dark
\end{CodeBlock}

A legacy entry resolves into one choice:

\begin{CodeBlock}
choice id: default
family id: default
mode id: undefined
label: Default
theme: default
\end{CodeBlock}

The existing frame model can continue storing \texttt{a\allowbreak{}c\allowbreak{}t\allowbreak{}i\allowbreak{}v\allowbreak{}e\allowbreak{}S\allowbreak{}h\allowbreak{}e\allowbreak{}l\allowbreak{}l\allowbreak{}T\allowbreak{}h\allowbreak{}e\allowbreak{}m\allowbreak{}e\allowbreak{}I\allowbreak{}d} as the choice id. Consumers that inspect \texttt{F\allowbreak{}r\allowbreak{}a\allowbreak{}m\allowbreak{}e\allowbreak{}S\allowbreak{}h\allowbreak{}e\allowbreak{}l\allowbreak{}l\allowbreak{}T\allowbreak{}h\allowbreak{}e\allowbreak{}m\allowbreak{}e\allowbreak{}C\allowbreak{}h\allowbreak{}a\allowbreak{}n\allowbreak{}g\allowbreak{}e} get the richer facts:

\begin{CodeBlock}
{
  shellTheme: {
    id: 'dogfood:dark',
    label: 'DOGFOOD / Dark',
    theme: dogfoodDark,
  },
  shellThemeSpec: dogfoodFamily,
  shellThemeId: 'dogfood',
  modeId: 'dark',
  modeLabel: 'Dark',
  ctx,
}
\end{CodeBlock}

\section{Settings UX}\label{record-181-settings-ux}

The first implementation keeps one settings row:

\begin{CodeBlock}
Shell theme  DOGFOOD / Dark
\end{CodeBlock}

Activating the row cycles the resolved choices:

% Presentation of original arrow notation. See flow-manifest.json.
\begin{center}
\begin{tikzpicture}[book flow node/.style={draw=BookHeading,line width=.65pt,fill=BookPanel,rounded corners=1.5pt,align=center,inner sep=5pt,minimum height=10mm,font=\fontsize{9}{11}\selectfont,execute at begin node={\hyphenpenalty=10000\exhyphenpenalty=10000}},book flow edge/.style={-{Latex[length=2mm]},draw=BookHeading,line width=.7pt}]
\node[book flow node,text width=96mm] (n0) at (0,0.0) {DOGFOOD /\allowbreak{} Dark};
\node[book flow node,text width=96mm] (n1) at (0,-1.55) {DOGFOOD /\allowbreak{} Light};
\node[book flow node,text width=96mm] (n2) at (0,-3.1) {Cabinet of Curiosities};
\node[book flow node,text width=96mm] (n3) at (0,-4.65) {.\allowbreak{}.\allowbreak{}.\allowbreak{}};
\draw[book flow edge] (n0) -- (n1);
\draw[book flow edge] (n1) -- (n2);
\draw[book flow edge] (n2) -- (n3);
\end{tikzpicture}
\end{center}


This is intentionally conservative. A later \#311/\#315 slice can split family and mode into separate controls once the app-frame has enough option metadata to render that well.

\section{Compatibility}\label{record-181-compatibility}

Compatibility is a hard requirement:

\begin{itemize}
\tightlist
\item
  Existing \texttt{F\allowbreak{}r\allowbreak{}a\allowbreak{}m\allowbreak{}e\allowbreak{}S\allowbreak{}h\allowbreak{}e\allowbreak{}l\allowbreak{}l\allowbreak{}T\allowbreak{}h\allowbreak{}e\allowbreak{}m\allowbreak{}e} entries with \texttt{t\allowbreak{}h\allowbreak{}e\allowbreak{}m\allowbreak{}e} still compile.
\item
  Existing apps that read \texttt{c\allowbreak{}h\allowbreak{}a\allowbreak{}n\allowbreak{}g\allowbreak{}e\allowbreak{}.\allowbreak{}s\allowbreak{}h\allowbreak{}e\allowbreak{}l\allowbreak{}l\allowbreak{}T\allowbreak{}h\allowbreak{}e\allowbreak{}m\allowbreak{}e\allowbreak{}.\allowbreak{}t\allowbreak{}h\allowbreak{}e\allowbreak{}m\allowbreak{}e} still receive a concrete selected shell theme.
\item
  Mode-aware apps can inspect \texttt{c\allowbreak{}h\allowbreak{}a\allowbreak{}n\allowbreak{}g\allowbreak{}e\allowbreak{}.\allowbreak{}s\allowbreak{}h\allowbreak{}e\allowbreak{}l\allowbreak{}l\allowbreak{}T\allowbreak{}h\allowbreak{}e\allowbreak{}m\allowbreak{}e\allowbreak{}S\allowbreak{}p\allowbreak{}e\allowbreak{}c}, \texttt{c\allowbreak{}h\allowbreak{}a\allowbreak{}n\allowbreak{}g\allowbreak{}e\allowbreak{}.\allowbreak{}s\allowbreak{}h\allowbreak{}e\allowbreak{}l\allowbreak{}l\allowbreak{}T\allowbreak{}h\allowbreak{}e\allowbreak{}m\allowbreak{}e\allowbreak{}I\allowbreak{}d}, and optional mode facts for family identity.
\item
  Existing app-frame tests for simple theme cycling keep passing.
\item
  \texttt{c\allowbreak{}l\allowbreak{}o\allowbreak{}n\allowbreak{}e\allowbreak{}S\allowbreak{}h\allowbreak{}e\allowbreak{}l\allowbreak{}l\allowbreak{}T\allowbreak{}h\allowbreak{}e\allowbreak{}m\allowbreak{}e\allowbreak{}C\allowbreak{}o\allowbreak{}n\allowbreak{}t\allowbreak{}e\allowbreak{}x\allowbreak{}t\allowbreak{}(\allowbreak{})} still receives one concrete resolved theme.
\end{itemize}

\section{Tests}\label{record-181-tests}

The slice is complete when tests prove:

\begin{enumerate}
\def\labelenumi{\arabic{enumi}.}
\tightlist
\item
  Mode-aware shell themes flatten into stable choices.
\item
  Settings cycling applies a mode choice and reports family/mode facts.
\item
  Old single-mode shell themes still cycle and apply.
\item
  DOGFOOD publishes one \texttt{d\allowbreak{}o\allowbreak{}g\allowbreak{}f\allowbreak{}o\allowbreak{}o\allowbreak{}d} family with \texttt{d\allowbreak{}a\allowbreak{}r\allowbreak{}k} and \texttt{l\allowbreak{}i\allowbreak{}g\allowbreak{}h\allowbreak{}t} modes.
\item
  Preview/resolution helpers can inspect a mode choice without mutating the ambient default context until the setting is applied.
\end{enumerate}

\section{Acceptance Criteria}\label{record-181-acceptance-criteria}

\begin{itemize}
\tightlist
\item[$\square$]
  \texttt{F\allowbreak{}r\allowbreak{}a\allowbreak{}m\allowbreak{}e\allowbreak{}S\allowbreak{}h\allowbreak{}e\allowbreak{}l\allowbreak{}l\allowbreak{}T\allowbreak{}h\allowbreak{}e\allowbreak{}m\allowbreak{}e\allowbreak{}S\allowbreak{}p\allowbreak{}e\allowbreak{}c} supports concrete \texttt{F\allowbreak{}r\allowbreak{}a\allowbreak{}m\allowbreak{}e\allowbreak{}S\allowbreak{}h\allowbreak{}e\allowbreak{}l\allowbreak{}l\allowbreak{}T\allowbreak{}h\allowbreak{}e\allowbreak{}m\allowbreak{}e} entries and mode-aware \texttt{F\allowbreak{}r\allowbreak{}a\allowbreak{}m\allowbreak{}e\allowbreak{}S\allowbreak{}h\allowbreak{}e\allowbreak{}l\allowbreak{}l\allowbreak{}T\allowbreak{}h\allowbreak{}e\allowbreak{}m\allowbreak{}e\allowbreak{}F\allowbreak{}a\allowbreak{}m\allowbreak{}i\allowbreak{}l\allowbreak{}y} entries.
\item[$\square$]
  Runtime choice ids distinguish family and mode.
\item[$\square$]
  \texttt{F\allowbreak{}r\allowbreak{}a\allowbreak{}m\allowbreak{}e\allowbreak{}S\allowbreak{}h\allowbreak{}e\allowbreak{}l\allowbreak{}l\allowbreak{}T\allowbreak{}h\allowbreak{}e\allowbreak{}m\allowbreak{}e\allowbreak{}C\allowbreak{}h\allowbreak{}a\allowbreak{}n\allowbreak{}g\allowbreak{}e} includes optional mode facts.
\item[$\square$]
  DOGFOOD presents one DogFood shell theme family with dark/light modes.
\item[$\square$]
  Legacy single-mode themes remain supported.
\item[$\square$]
  Tests cover mode switching and context isolation.
\end{itemize}

\section{Slice Plan}\label{record-181-slice-plan}

\begin{enumerate}
\def\labelenumi{\arabic{enumi}.}
\tightlist
\item
  Add the design doc and tracker updates.
\item
  Add app-frame mode-aware shell theme types and choice resolution tests.
\item
  Update settings cycling and shell change notifications.
\item
  Convert DOGFOOD shell themes to one mode-aware family.
\item
  Update Theme Inspector, Theme Lab, docs tests, and changelog.
\item
  Run self-review and validation before opening the PR.
\end{enumerate}

